\chapter{有界自伴算子的谱定理}
\label{chap:bounded-spectral-theorem}
\index{谱定理}
\index{投影值测度}

本章给出有界自伴算子的谱定理.我们先建立有界算子的若干初等性质,然后给出谱定理的两种形式:第一种形式使用投影值测度,第二种形式使用直积分.前一种形式最适合函数演算和测量概率,后一种形式则更接近“广义特征向量”的结构.

\section{有界算子的初等性质}
\label{sec:bounded-operator-properties}

记 \(\mathcal{B}(\mathcal{H})\) 为 Hilbert 空间 \(\mathcal{H}\) 上所有有界线性算子组成的 Banach 空间.对任意 \(A\in\mathcal{B}(\mathcal{H})\),附录~\ref{app:subsec:hilbert} 回顾过存在唯一的算子 \(A^{*}\in\mathcal{B}(\mathcal{H})\),称为 \(A\) 的伴随,满足
\[
\langle\phi,A\psi\rangle=\langle A^{*}\phi,\psi\rangle
\]
对所有 \(\phi,\psi\in\mathcal{H}\) 成立.若 \(A^{*}=A\),则称 \(A\) 为自伴算子.若
\[
\langle\psi,A\psi\rangle\geq 0
\label{eq:bounded-nonnegative}
\]
对所有 \(\psi\in\mathcal{H}\) 成立,则称 \(A\) 为非负算子.

\begin{proposition}[伴随算子的范数]\label{prop:adjoint-norm}
对任意 \(A\in\mathcal{B}(\mathcal{H})\),有
\[
\|A^{*}\|=\|A\|,
\qquad
\|A^{*}A\|=\|A\|^{2}.
\]
特别地,若 \(A\) 自伴,则
\[
\|A^{2}\|=\|A\|^{2}.
\]
\end{proposition}

\begin{proof}
算子范数也可以写成
\[
\|A\|=\sup_{\|\psi\|=1}\|A\psi\|.
\]
另一方面,对任意 \(\phi\in\mathcal{H}\),由 Cauchy--Schwarz 不等式以及取 \(\chi\) 为 \(\phi\) 的适当倍数可得
\[
\|\phi\|=\sup_{\|\chi\|=1}|\langle\chi,\phi\rangle|.
\]
因此
\[
\|A\|=\sup_{\|\phi\|=\|\psi\|=1}|\langle\phi,A\psi\rangle|.
\]
由此得到
\begin{align*}
\|A^{*}\|
&=\sup_{\|\phi\|=\|\psi\|=1}|\langle\phi,A^{*}\psi\rangle|\\
&=\sup_{\|\phi\|=\|\psi\|=1}|\langle A\phi,\psi\rangle|\\
&=\sup_{\|\phi\|=\|\psi\|=1}|\langle\psi,A\phi\rangle|=\|A\|.
\end{align*}
同时,\(\|A^{*}A\|\leq\|A^{*}\|\,\|A\|=\|A\|^{2}\).反过来,
\begin{align*}
\|A^{*}A\|
&=\sup_{\|\phi\|=\|\psi\|=1}|\langle\phi,A^{*}A\psi\rangle|\\
&=\sup_{\|\phi\|=\|\psi\|=1}|\langle A\phi,A\psi\rangle|\\
&\geq\sup_{\|\psi\|=1}|\langle A\psi,A\psi\rangle|=\|A\|^{2}.
\end{align*}
故第二个等式也成立.若 \(A\) 自伴,则 \(A^{*}A=A^{2}\).
\end{proof}

\def\ChapterSevenLater{
\section{有界自伴算子的谱定理(二)}
\label{sec:spectral-theorem-second}

正如第~\ref{sec:multiplication-operators} 节已经指出的,谱定理的一种表述是:每个自伴算子都与一个乘法算子酉等价.对可分 Hilbert 空间上的有界自伴算子 \(A\),这意味着 \(A\) 与 \(L^{2}(X,\mu)\) 上的乘法算子 \(M_{h}\) 酉等价,其中 \((X,\mu)\) 是 \(\sigma\)-有限测度空间,\(h\) 是可测实值函数,且
\[
(M_{h}\psi)(\lambda)=h(\lambda)\psi(\lambda).
\]

虽然乘法算子形式(定理~\ref{thm:spectral-theorem-multiplication})很容易陈述,但还有一种更好的形式,它使用直积分的概念.把取值于 Hilbert 空间 \(\mathcal{H}\) 的 \(L^{2}\) 空间推广到更一般的情形并不困难;在直积分中,我们再向前一步,允许 Hilbert 空间随点而变化.我们从测度空间 \((X,\mu)\) 开始,对每个 \(\lambda\in X\) 给定一个 Hilbert 空间 \(\mathcal{H}_{\lambda}\).直积分中的元素是函数 \(s\),其在每个点的取值 \(s(\lambda)\) 属于 \(\mathcal{H}_{\lambda}\).给定 \(X\) 上的实值可测函数 \(h\),就可以用 \(h\) 乘直积分中的元素 \(s\).

谱定理的直积分形式说,有界自伴算子 \(A\) 与直积分上的乘法算子酉等价.与普通 \(L^{2}\) 空间上的乘法算子相比,直积分有几项优点.第一,集合 \(X\) 与函数 \(h\) 变成规范的:\(X\) 就是 \(A\) 的谱,而 \(h(\lambda)=\lambda\).第二,直积分自带“广义特征向量”的概念,因为 \(\mathcal{H}_{\lambda}\) 可以看成特征值为 \(\lambda\) 的广义特征向量空间.(一般来说,\(\mathcal{H}_{\lambda}\) 并不包含在直积分 Hilbert 空间中;因此,直积分为“不属于算子作用的 Hilbert 空间的特征向量”这一想法提供了严格含义.)第三,直积分给出了按酉等价分类自伴算子的简单方法:两个自伴算子酉等价,当且仅当它们的直积分表示在自然意义下等价(见命题~\ref{prop:bounded-spectral-classification}).

如果只希望得到普通乘法算子版本的简洁结论,那么使用与直积分证明完全相同的方法即可证明它(见定理~\ref{thm:spectral-theorem-multiplication}).不过,对单个自伴算子而言,直积分版本可以说是最完整的谱定理表述.

\subsection{直积分}
\label{subsec:direct-integral}

现在定义直积分.设 \(\mu\) 是集合 \(X\) 上某个 \(\sigma\)-代数 \(\Omega\) 的 \(\sigma\)-有限测度.假设对每个 \(\lambda\in X\),都有一个带内积 \(\langle\cdot,\cdot\rangle_{\lambda}\) 的可分 Hilbert 空间 \(\mathcal{H}_{\lambda}\).我们希望定义关于 \(\mu\) 的 \(\mathcal{H}_{\lambda}\) 的直积分.直积分中的元素将是截面,即函数 \(s\),其值属于这些 Hilbert 空间的并,并满足
\[
s(\lambda)\in\mathcal{H}_{\lambda}
\]
对每个 \(\lambda\in X\) 成立.若右端积分有限,我们希望用
\[
\|s\|^{2}=\int_{X}\langle s(\lambda),s(\lambda)\rangle_{\lambda}\,\mathrm{d}\mu(\lambda)
\]
定义截面 \(s\) 的范数;两个有限范数截面 \(s_{1},s_{2}\) 的内积则应为
\[
\langle s_{1},s_{2}\rangle
:=\int_{X}\langle s_{1}(\lambda),s_{2}(\lambda)\rangle_{\lambda}\,\mathrm{d}\mu(\lambda).
\]

这种描述的问题在于我们还没有说明可测性.由于 \(s(\lambda)\) 对每个 \(\lambda\) 属于不同的空间,直接询问截面 \(s\) 是否可测并没有意义.因此,需要引入额外结构,从而产生可测性的概念.(第一次阅读时,可以把可测性问题看作一个技术细节.)

处理可测性的一种方法,是为每个 \(\mathcal{H}_{\lambda}\) 选择同时的标准正交基.为允许不同空间具有不同维数,我们稍微修改标准正交基的概念:若向量族 \(\{e_{j}\}\) 满足 \(\langle e_{j},e_{k}\rangle=0\)(\(j\neq k\)),每个 \(e_{j}\) 的范数为 0 或 1,且它们张成空间的闭包是整个 \(\mathcal{H}\),就称它为 Hilbert 空间 \(\mathcal{H}\) 的标准正交基.也就是说,我们允许基中的某些向量为零,而非零向量组成通常意义下的标准正交基.

对于 Hilbert 空间族 \(\{\mathcal{H}_{\lambda}\}\),同时标准正交基是截面族 \(\{e_{j}(\cdot)\}_{j=1}^{\infty}\),使得对每个 \(\lambda\),\(\{e_{j}(\lambda)\}_{j=1}^{\infty}\) 是 \(\mathcal{H}_{\lambda}\) 的上述意义下的标准正交基.若函数 \(\lambda\mapsto\dim\mathcal{H}_{\lambda}\) 从 \(X\) 到 \([0,\infty]\) 可测,则可以选择同时标准正交基,使 \(\langle e_{j}(\lambda),e_{k}(\lambda)\rangle_{\lambda}\) 对所有 \(j,k\) 都可测.选定这样的基后,若对每个 \(j\),函数
\[
\lambda\longmapsto\langle e_{j}(\lambda),s(\lambda)\rangle_{\lambda}
\]
都是复值可测函数,就称截面 \(s\) 可测.对基本身的假设意味着每个 \(e_{j}\) 都是可测截面.

把满足上述内积可测条件的同时标准正交基的选择称为 Hilbert 空间族 \(\{\mathcal{H}_{\lambda}\}\) 的可测性结构.若 \(s_{1},s_{2}\) 是两个可测截面,则函数
\[
\lambda\longmapsto\langle s_{1}(\lambda),s_{2}(\lambda)\rangle_{\lambda}
=\sum_{j=1}^{\infty}
\langle s_{1}(\lambda),e_{j}(\lambda)\rangle_{\lambda}
\langle e_{j}(\lambda),s_{2}(\lambda)\rangle_{\lambda}
\]
也是可测函数.

\begin{definition}[Hilbert 空间的直积分]\label{def:direct-integral}
给定以下结构:
\begin{enumerate}
\item \(\sigma\)-有限测度空间 \((X,\Omega,\mu)\);
\item 一个可分 Hilbert 空间族 \(\{\mathcal{H}_{\lambda}\}_{\lambda\in X}\),其维数函数可测;
\item 该 Hilbert 空间族上的可测性结构.
\end{enumerate}
则关于 \(\mu\) 的直积分记作
\[
\int_{X}^{\oplus}\mathcal{H}_{\lambda}\,\mathrm{d}\mu(\lambda),
\]
它是所有满足以下条件的可测截面按几乎处处相等取等价类所得的空间:
\[
\|s\|^{2}:=\int_{X}\langle s(\lambda),s(\lambda)\rangle_{\lambda}\,\mathrm{d}\mu(\lambda)<\infty.
\]
两个截面 \(s_{1},s_{2}\) 的内积定义为
\[
\langle s_{1},s_{2}\rangle
:=\int_{X}\langle s_{1}(\lambda),s_{2}(\lambda)\rangle_{\lambda}\,\mathrm{d}\mu(\lambda).
\]
\end{definition}

为什么两个有限范数截面的内积积分有限？注意
\[
|\langle s_{1}(\lambda),s_{2}(\lambda)\rangle_{\lambda}|
\leq\|s_{1}(\lambda)\|_{\lambda}\|s_{2}(\lambda)\|_{\lambda}.
\]
按假设,\(\|s_{j}(\lambda)\|_{\lambda}\)(\(j=1,2\))是 \(\lambda\) 的平方可积函数,而两个平方可积函数的乘积可积.因此,内积定义中的被积函数也可积.类似于证明 \(L^{2}\) 空间完备性的论证,可以证明 Hilbert 空间的直积分仍是 Hilbert 空间.

直积分构造有两个重要特例.第一,如果每个 \(\mathcal{H}_{\lambda}\) 都只是 \(\mathbb{C}\),那么直积分(带有显然的可测性结构)就是 \(L^{2}(X,\mu)\).第二,若 \(X=\{\lambda_{1},\lambda_{2},\ldots\}\) 是可数集,\(\Omega\) 是 \(X\) 的所有子集组成的 \(\sigma\)-代数,\(\mu\) 是计数测度,则直积分就是 Hilbert 空间直和(定义~\ref{app:def:hilbert-direct-sum}).

给定一个直积分,假设某个 \(\lambda_{0}\in X\) 的单点集 \(\{\lambda_{0}\}\) 可测,且 \(c:=\mu(\{\lambda_{0}\})>0\).则可以把 \(\mathcal{H}_{\lambda_{0}}\) 等距嵌入直积分:将 \(\psi\in\mathcal{H}_{\lambda_{0}}\) 映到截面
\[
s(\lambda)=
\begin{cases}
c^{-1/2}\psi,&\lambda=\lambda_{0},\\
0,&\lambda\neq\lambda_{0}.
\end{cases}
\]
即使 \(\mu(\{\lambda_{0}\})=0\),我们仍可把 \(\mathcal{H}_{\lambda_{0}}\) 看成直积分的某种“广义子空间”.

\begin{theorem}[谱定理:第二种形式]\label{thm:spectral-theorem-direct-integral}
若 \(A\in\mathcal{B}(\mathcal{H})\) 自伴,则在 \(\sigma(A)\) 上存在 \(\sigma\)-有限测度 \(\mu\)、直积分
\[
\int_{\sigma(A)}^{\oplus}\mathcal{H}_{\lambda}\,\mathrm{d}\mu(\lambda),
\]
以及从 \(\mathcal{H}\) 到该直积分的酉映射 \(U\),使得对直积分中的所有截面 \(s\),
\[
\bigl[UAU^{-1}s\bigr](\lambda)=\lambda s(\lambda).
\label{eq:spectral-theorem-direct-integral}
\]
\end{theorem}

定理~\ref{thm:spectral-theorem-direct-integral} 的证明将在下一章与第一种谱定理的证明一同给出.现在先解释它的含义.可以把酉映射 \(U\) 看成将原 Hilbert 空间 \(\mathcal{H}\) 识别为关于 \(A\) 的谱的某个直积分.在这种识别下,自伴算子 \(A\) 变成乘以 \(\lambda\) 的算子,即把截面 \(s(\lambda)\) 送到 \(\lambda s(\lambda)\) 的映射.粗略地说,在这个识别下,\(A\) 在每个空间 \(\mathcal{H}_{\lambda}\) 上都像 \(\lambda I\) 那样作用.因此,可以把每个 \(\mathcal{H}_{\lambda}\) 看成 \(A\) 对应于谱点 \(\lambda\) 的“特征空间”.当然,除非 \(\mu(\{\lambda\})>0\),Hilbert 空间 \(\mathcal{H}_{\lambda}\) 并不真正包含在 \(\mathcal{H}\) 中;不过仍可把其中元素视为 \(A\) 的“广义特征向量”.

直积分形式很容易导出有界自伴算子的分类结果(见本节后面的命题~\ref{prop:bounded-spectral-classification}).与此同时,正如前面提到的,定理~\ref{thm:spectral-theorem-direct-integral} 的证明方法也会给出普通 \(L^{2}\) 空间上乘法算子形式的谱定理.

\begin{theorem}[谱定理:乘法算子形式]\label{thm:spectral-theorem-multiplication}
设 \(A\in\mathcal{B}(\mathcal{H})\) 自伴.则存在 \(\sigma\)-有限测度空间 \((X,\mu)\)、有界可测实值函数 \(h\) 以及酉映射 \(U:\mathcal{H}\to L^{2}(X,\mu)\),使得对所有 \(\psi\in L^{2}(X,\mu)\),
\[
\bigl[UAU^{-1}(\psi)\bigr](\lambda)=h(\lambda)\psi(\lambda).
\]
\end{theorem}

回到谱定理的直积分表述.它给出了函数演算的简单描述.

\begin{proposition}[直积分中的函数演算]\label{prop:direct-integral-functional-calculus}
设 \(A\in\mathcal{B}(\mathcal{H})\) 自伴,\(U\) 是定理~\ref{thm:spectral-theorem-direct-integral} 中的酉映射.则对 \(\sigma(A)\) 上任意有界可测函数 \(f\),有
\[
\bigl[Uf(A)U^{-1}(s)\bigr](\lambda)=f(\lambda)s(\lambda).
\label{eq:direct-integral-functional-calculus}
\]
\end{proposition}

因此,粗略地说,\(f(A)\) 在每个“广义特征空间” \(\mathcal{H}_{\lambda}\) 上等于 \(f(\lambda)I\).当 \(f\) 是多项式时,命题~\ref{prop:direct-integral-functional-calculus} 直接由式~\eqref{eq:spectral-theorem-direct-integral} 得到;连续函数的情形由一致极限得到;一般函数则可用第~8 章的极限论证,尤其是其中的练习 3,得到.

下面讨论第二种谱定理应当具有怎样的唯一性.存在一个“平凡”的非唯一性来源:某些 \(\mathcal{H}_{\lambda}\) 可能是零空间.令 \(E_{0}\) 为满足 \(\dim\mathcal{H}_{\lambda}=0\) 的 \(\lambda\) 的集合.即使 \(\mu(E_{0})>0\),\(E_{0}\) 也不会对截面的范数产生贡献,因为每个截面在 \(E_{0}\) 上自动为零.因此可以定义新测度
\[
\widetilde{\mu}(E)=\mu(E\cap E_{0}^{c}),
\]
它在 \(E_{0}^{c}\) 上与 \(\mu\) 一致,在 \(E_{0}\) 上为零.关于 \(\mu\) 和关于 \(\widetilde{\mu}\) 的直积分是“无法区分的”.所以总可以修改直积分,使 \(\mu\)-几乎处处都有 \(\dim\mathcal{H}_{\lambda}>0\).

另一方面,与定理~\ref{thm:spectral-theorem-first} 中的投影值测度 \(\mu_{A}\) 不同,定理~\ref{thm:spectral-theorem-direct-integral} 中的普通测度 \(\mu\) 并非唯一,而只在等价意义下唯一.这里两个给定可测空间上的 \(\sigma\)-有限测度等价,是指它们有完全相同的零测集.给定 \(\mu\) 后,Hilbert 空间 \(\mathcal{H}_{\lambda}\) 只在酉等价意义下唯一,也就是说,唯一确定的只是各空间的维数.甚至 \(\dim\mathcal{H}_{\lambda}\) 也只在一个 \(\mu\)-零集之外唯一.事实证明,本段及上一段中所述的非唯一性来源就是全部非唯一性来源.

\begin{proposition}[定理~\ref{thm:spectral-theorem-direct-integral} 的唯一性]\label{prop:direct-integral-uniqueness}
设 \(A\in\mathcal{B}(\mathcal{H})\) 自伴,考虑定理~\ref{thm:spectral-theorem-direct-integral} 的两个不同直积分表示:第一个使用测度 \(\mu^{(1)}\) 和 Hilbert 空间 \(\mathcal{H}_{\lambda}^{(1)}\),第二个使用测度 \(\mu^{(2)}\) 和 Hilbert 空间 \(\mathcal{H}_{\lambda}^{(2)}\).若对 \(j=1,2\),\(\dim\mathcal{H}_{\lambda}^{(j)}>0\) 对 \(\mu^{(j)}\)-几乎处处的 \(\lambda\) 成立,则 \(\mu^{(1)}\) 与 \(\mu^{(2)}\) 互相绝对连续,并且
\[
\dim\mathcal{H}_{\lambda}^{(1)}=\dim\mathcal{H}_{\lambda}^{(2)}
\]
对 \(\mu^{(j)}\)-几乎处处的 \(\lambda\) 成立(\(j=1,2\)).
\end{proposition}

本唯一性结果的证明概要见下一章末尾.

应把定理~\ref{thm:spectral-theorem-direct-integral} 看成此前形式(定理~\ref{thm:spectral-theorem-first})的细化.事实上,给定可测集 \(E\subset\sigma(A)\),令 \(V_{E}\) 为所有支集在 \(E\) 上的截面(即在 \(E^{c}\) 上 \(\mu\)-几乎处处为零的截面)的等价类组成的空间.显然这是一个闭子空间.令 \(P_{E}\) 为投影到 \(V_{E}\) 上的正交投影,并定义
\[
\mu_{A}(E)=U^{-1}P_{E}U.
\label{eq:direct-integral-pvm}
\]
不难验证,\(\mu_{A}\) 是以 \(\mathcal{B}(\mathcal{H})\) 中算子为值的投影值测度,并且
\[
\int_{\sigma(A)}\lambda\,\mathrm{d}\mu_{A}(\lambda)=A.
\]

注意,关于 \(A\) 的两个谱定理版本都涉及测度:第一种是投影值测度 \(\mu_{A}\),第二种是取非负实数值的普通测度 \(\mu\).下面的结果说明这两个测度之间的关系.

\begin{proposition}[两种谱测度的零集]\label{prop:spectral-measures-equivalence}
设 \(A\in\mathcal{B}(\mathcal{H})\) 自伴,\(\mu_{A}\) 是定理~\ref{thm:spectral-theorem-first} 给出的投影值测度,\(\mu\) 是定理~\ref{thm:spectral-theorem-direct-integral} 中的实值测度.若 \(\dim\mathcal{H}_{\lambda}>0\) 对 \(\mu\)-几乎处处成立,则对任意 Borel 集 \(E\subset\sigma(A)\),
\[
\mu_{A}(E)=0\quad\Longleftrightarrow\quad\mu(E)=0.
\]
\end{proposition}

这里 \(\mu_{A}(E)=0\) 中的 0 是零算子,而 \(\mu(E)=0\) 中的 0 是实数.不过可以把命题~\ref{prop:spectral-measures-equivalence} 理解为说,\(\mu_{A}\) 与 \(\mu\) 在通常的测度论意义下等价,即具有完全相同的零测集.

\begin{proof}
如前所述,给定定理~\ref{thm:spectral-theorem-direct-integral} 中的直积分,可以通过式~\eqref{eq:direct-integral-pvm} 构造投影值测度;该投影值测度满足
\[
\int_{\sigma(A)}\lambda\,\mathrm{d}\mu_{A}(\lambda)=A.
\]
由定理~\ref{thm:spectral-theorem-first} 的唯一性,它必须就是该定理中的 \(\mu_{A}\).

若 \(\mu(E)=0\),则支集在 \(E\) 上的任意截面几乎处处为零,因而表示直积分中的零元素.此时 \(V_{E}=\{0\}\),由式~\eqref{eq:direct-integral-pvm} 得 \(\mu_{A}(E)=0\).反之,若 \(\mu(E)>0\),由于 \(\mu\) 是 \(\sigma\)-有限的,\(E\) 含有可测子集 \(F\),满足 \(0<\mu(F)<\infty\).令 \(\{e_{j}(\cdot)\}\) 为直积分的可测性结构,并定义
\[
s(\lambda)=
\begin{cases}
\displaystyle\sum_{j=1}^{\infty}2^{-j}e_{j}(\lambda),&\lambda\in F,\\
0,&\lambda\in F^{c}.
\end{cases}
\]
则
\[
\langle s(\lambda),e_{j}(\lambda)\rangle_{\lambda}
=2^{-j}\langle e_{j}(\lambda),e_{j}(\lambda)\rangle_{\lambda}\mathbf{1}_{F}(\lambda),
\]
对每个 \(j\) 都是可测函数,所以 \(s\) 可测.由 \(\mathcal{H}_{\lambda}\) 在 \(\mu\)-几乎处处非零的假设,\(s\) 在 \(F\) 上几乎处处非零,因而范数为正.又因 \(\|s(\lambda)\|\leq1\),且 \(F\) 的测度有限,\(s\) 的范数有限.因此 \(V_{E}\neq\{0\}\),从而 \(\mu_{A}(E)\neq0\).
\end{proof}

若 \(A_{1},A_{2}\) 分别作用在 Hilbert 空间 \(\mathcal{H}_{1},\mathcal{H}_{2}\) 上,若存在酉映射 \(U:\mathcal{H}_{1}\to\mathcal{H}_{2}\),使得
\[
A_{2}=UA_{1}U^{-1},
\]
则称 \(A_{1},A_{2}\) 酉等价.利用命题~\ref{prop:direct-integral-uniqueness},可以给出可分 Hilbert 空间上有界自伴算子按酉等价的分类.对给定的有界自伴算子 \(A\),称函数
\[
\lambda\longmapsto\dim\mathcal{H}_{\lambda}
\]
为 \(A\) 的重数函数.它与直积分分解的选择无关,除了可能在一个零测集上不同.事实证明,有界自伴算子在酉等价意义下由三项数据刻画:作为集合的谱 \(\sigma(A)\)、定理~\ref{thm:spectral-theorem-direct-integral} 中测度 \(\mu\) 的等价类,以及重数函数.

\begin{proposition}[有界自伴算子的分类]\label{prop:bounded-spectral-classification}
设 \(A_{1},A_{2}\) 分别是可分 Hilbert 空间 \(\mathcal{H}_{1},\mathcal{H}_{2}\) 上的有界自伴算子.为它们分别选择定理~\ref{thm:spectral-theorem-direct-integral} 中的直积分表示,并要求 \(\dim\mathcal{H}_{\lambda}>0\) 对相应测度 \(\mu_{j}\)-几乎处处成立(\(j=1,2\)).则 \(A_{1},A_{2}\) 酉等价,当且仅当:
\begin{enumerate}
\item \(\sigma(A_{1})=\sigma(A_{2})\);
\item \(\mu_{1},\mu_{2}\) 互相绝对连续;
\item \(A_{1},A_{2}\) 的重数函数除一个零测集外相同.
\end{enumerate}
\end{proposition}

证明见练习~\ref{ex:prove-spectral-classification}.

\section{练习}
\label{sec:chapter7-exercises}

\begin{exercise}[同时特征向量]\label{ex:simultaneous-eigenvector}
设 \(A,B\) 是作用在非零有限维向量空间上的可交换线性算子.
\begin{enumerate}[label=(\alph*)]
\item 证明 \(A\) 的每个特征子空间在 \(B\) 下不变;
\item 证明 \(A,B\) 至少有一个同时特征向量,即存在非零向量 \(v\) 及常数 \(\lambda,\mu\in\mathbb{C}\),使得 \(Av=\lambda v\),\(Bv=\mu v\).
\end{enumerate}
\end{exercise}

\begin{exercise}[正规算子的特征向量]\label{ex:normal-eigenvector-adjoint}
设 \(A\in\mathcal{B}(\mathcal{H})\) 是正规算子,即 \(AA^{*}=A^{*}A\).若对某个 \(\psi\in\mathcal{H}\) 和 \(\lambda\in\mathbb{C}\),有 \(A\psi=\lambda\psi\),证明 \(A^{*}\psi=\overline{\lambda}\psi\).

\emph{提示:}计算 \(\|(A^{*}-\overline{\lambda}I)\psi\|\).
\end{exercise}

\begin{exercise}[不变子空间的正交补]\label{ex:invariant-orthogonal-complement}
设 \(V\) 是 \(\mathcal{H}\) 的闭子空间,并且在有界算子 \(A\) 下不变,即对所有 \(\psi\in V\),都有 \(A\psi\in V\).证明 \(V^{\perp}\) 在 \(A^{*}\) 下不变.
\end{exercise}

\begin{exercise}[有限维正规算子的标准正交基]\label{ex:finite-normal-orthonormal-basis}
\begin{enumerate}[label=(\alph*)]
\item 设 \(\mathcal{H}\) 是复有限维 Hilbert 空间,\(A\) 是练习~\ref{ex:normal-eigenvector-adjoint} 意义下的正规算子.证明存在由 \(A\) 与 \(A^{*}\) 的同时特征向量组成的 \(\mathcal{H}\) 的标准正交基.
\emph{提示:}使用练习~\ref{ex:simultaneous-eigenvector} 和练习~\ref{ex:invariant-orthogonal-complement}.
\item 设 \(A\) 是复有限维 Hilbert 空间上的线性算子,并假设存在由 \(A\) 的特征向量组成的标准正交基.证明 \(A\) 与 \(A^{*}\) 相交换.
\end{enumerate}
\end{exercise}

\begin{exercise}[逆算子的伴随]\label{ex:adjoint-inverse}
设 \(A\in\mathcal{B}(\mathcal{H})\) 在 \(\mathcal{B}(\mathcal{H})\) 中有逆 \(A^{-1}\).证明
\[
(A^{-1})^{*}A^{*}=A^{*}(A^{-1})^{*}=I.
\]
推出 \(A^{*}\) 可逆,且 \((A^{*})^{-1}=(A^{-1})^{*}\).
\end{exercise}

\begin{exercise}[酉算子的谱]\label{ex:unitary-spectrum}
设 \(U\) 是 \(\mathcal{H}\) 上的酉算子(定义~\ref{app:def:unitary}).证明 \(U\) 的谱包含于单位圆.

\emph{提示:}将 \(U-\lambda I\) 写成 \((-\lambda)(I-U/\lambda)\) 或 \(U(I-\lambda U^{-1})\),证明所有满足 \(|\lambda|\neq1\) 的 \(\lambda\) 都属于预解集.(原文提示中的算子应为 \(U\),而不是 \(\lambda\).)
\end{exercise}

\begin{exercise}[非负算子的谱]\label{ex:nonnegative-spectrum}
设 \(A\in\mathcal{B}(\mathcal{H})\) 自伴且非负,即满足式~\eqref{eq:bounded-nonnegative}.证明 \(A\) 的谱包含于区间 \([0,\infty)\).

\begin{note}
反过来,若 \(A\) 自伴且 \(\sigma(A)\subset[0,\infty)\),则 \(A\) 非负;见第~8 章练习 2.
\end{note}
\end{exercise}

\begin{exercise}[可逆性的稳定性]\label{ex:invertibility-stability}
设 \(A\in\mathcal{B}(\mathcal{H})\) 可逆.证明存在 \(\varepsilon>0\),使得对满足 \(\|B-A\|<\varepsilon\) 的所有 \(B\in\mathcal{B}(\mathcal{H})\),\(B\) 也可逆.

\emph{提示:}使用与命题~\ref{prop:spectrum-basic} 证明中相同的幂级数方法.
\end{exercise}

\begin{exercise}[自伴算子预解算子的范数]\label{ex:selfadjoint-resolvent-norm}
假设 \(A\in\mathcal{B}(\mathcal{H})\) 自伴.
\begin{enumerate}[label=(\alph*)]
\item 若 \(\lambda_{0}\in\mathbb{C}\) 属于 \(A\) 的预解集,证明
\[
\|(A-\lambda_{0}I)^{-1}\|
=\frac{1}{d(\lambda_{0},\sigma(A))},
\qquad
d(\lambda_{0},\sigma(A)):=\inf_{\lambda\in\sigma(A)}|\lambda-\lambda_{0}|.
\]
\emph{提示:}把 \((A-\lambda_{0}I)^{-1}\) 看成 \(A\) 的函数演算中的函数.
\item 给定 \(\lambda_{0}\in\mathbb{C}\),假设存在非零 \(\psi\in\mathcal{H}\),使得
\[
\|A\psi-\lambda_{0}\psi\|<\varepsilon\|\psi\|.
\]
证明存在 \(\lambda\in\sigma(A)\),使得 \(|\lambda-\lambda_{0}|<\varepsilon\).
\end{enumerate}
\end{exercise}

\begin{exercise}[正交投影的乘积]\label{ex:pvm-orthogonality}
设 \(V_{1},V_{2}\) 是 \(\mathcal{H}\) 的两个闭子空间,\(P_{1},P_{2}\) 是相应的正交投影.证明 \(V_{1}\perp V_{2}\) 当且仅当 \(P_{1}P_{2}=0\).
\end{exercise}

\begin{exercise}[投影值测度的交集]\label{ex:pvm-intersection}
设 \(\mu\) 是 \((X,\Omega)\) 上的投影值测度.证明对任意 \(E_{1},E_{2}\in\Omega\),\(\mu(E_{1})\mu(E_{2})\) 是投影到闭子空间
\[
\operatorname{Range}(\mu(E_{1}))\cap\operatorname{Range}(\mu(E_{2}))
\]
上的投影.

\emph{提示:}将 \(E_{1}\) 写为 \(E_{1}=(E_{1}\cap E_{2})\cup(E_{1}\setminus E_{2})\),并使用练习~\ref{ex:pvm-orthogonality}.
\end{exercise}

\begin{exercise}[证明谱分类定理]\label{ex:prove-spectral-classification}
证明命题~\ref{prop:bounded-spectral-classification}.

\emph{提示:}使用命题~\ref{prop:direct-integral-uniqueness} 和 Radon--Nikodym 定理(定理~\ref{app:thm:radon-nikodym}).
\end{exercise}

}

\begin{proposition}[值域与伴随的核]\label{prop:range-kernel-adjoint}
对任意 \(A\in\mathcal{B}(\mathcal{H})\),有
\[
\bigl[\operatorname{Range}(A)\bigr]^{\perp}=\ker(A^{*}).
\]
\end{proposition}

\begin{proof}
若 \(\psi\in[\operatorname{Range}(A)]^{\perp}\),则对所有 \(\phi\in\mathcal{H}\),
\[
0=\langle\psi,A\phi\rangle=\langle A^{*}\psi,\phi\rangle.
\label{eq:range-kernel-proof}
\]
于是 \(A^{*}\psi=0\),即 \(\psi\in\ker(A^{*})\).反过来,若 \(\psi\in\ker(A^{*})\),则式~\eqref{eq:range-kernel-proof} 从右向左读出,说明 \(\psi\) 与每个形如 \(A\phi\) 的向量正交.因此 \(\psi\in[\operatorname{Range}(A)]^{\perp}\).
\end{proof}

下面定义有界算子的谱;它在无限维情形中扮演有限维线性代数中特征值集合的角色.

\begin{definition}[谱与预解集]\label{def:spectrum-resolvent}
对 \(A\in\mathcal{B}(\mathcal{H})\),若 \(\lambda\in\mathbb{C}\) 使得算子 \(A-\lambda I\) 有有界逆,则称 \(\lambda\) 属于 \(A\) 的预解集,记作 \(\rho(A)\).\(A\) 的谱 \(\sigma(A)\) 定义为预解集在 \(\mathbb{C}\) 中的补集.对预解集中的 \(\lambda\),算子 \((A-\lambda I)^{-1}\) 称为 \(\lambda\) 处的预解算子.
\end{definition}

这里说 \(A-\lambda I\) 有有界逆,是指存在有界算子 \(B\),使得
\[
(A-\lambda I)B=B(A-\lambda I)=I.
\]
如果 \(A\) 有界,且 \(A-\lambda I\) 是一一映射并将 \(\mathcal{H}\) 映到整个 \(\mathcal{H}\),那么由闭图定理(定理~\ref{app:thm:closed-graph})可知其逆映射必然有界.因此,\(\rho(A)\) 也可以描述为所有使 \(A-\lambda I\) 既一一又满射的 \(\lambda\) 所组成的集合.

\begin{proposition}[有界算子的谱]\label{prop:spectrum-basic}
对任意 \(A\in\mathcal{B}(\mathcal{H})\),有:
\begin{enumerate}
\item \(\sigma(A)\) 是 \(\mathbb{C}\) 中非空、闭且有界的集合;
\item 若 \(|\lambda|>\|A\|\),则 \(\lambda\in\rho(A)\).
\end{enumerate}
\end{proposition}

\begin{lemma}[Neumann 级数]\label{lem:neumann-series}
若 \(X\in\mathcal{B}(\mathcal{H})\) 满足 \(\|X\|<1\),则 \(I-X\) 可逆,且其逆由 \(\mathcal{B}(\mathcal{H})\) 中收敛的级数给出:
\[
(I-X)^{-1}=I+X+X^{2}+X^{3}+\cdots.
\label{eq:neumann-series}
\]
\end{lemma}

\begin{proof}
由算子范数的次乘性,有 \(\|X^{m}\|\leq\|X\|^{m}\).因此式~\eqref{eq:neumann-series} 右端的几何级数绝对收敛,从而在 Banach 空间 \(\mathcal{B}(\mathcal{H})\) 中收敛(见附录~\ref{app:subsec:banach}).在级数两侧乘以 \(I-X\) 后,除 \(I\) 外的各项全部相消,故其和就是 \(I-X\) 的逆.
\end{proof}

\begin{proof}[命题~\ref{prop:spectrum-basic} 的证明]
对任意非零 \(\lambda\in\mathbb{C}\),写成
\[
A-\lambda I=-\lambda\left(I-\frac{A}{\lambda}\right).
\]
若 \(|\lambda|>\|A\|\),则 \(\|A/\lambda\|<1\),由引理~\ref{lem:neumann-series},\(I-A/\lambda\) 可逆,从而 \(A-\lambda I\) 可逆,且
\[
(A-\lambda I)^{-1}
=-\frac{1}{\lambda}\left(I+\frac{A}{\lambda}+\frac{A^{2}}{\lambda^{2}}+\cdots\right).
\label{eq:resolvent-neumann}
\]
这证明了第二点,并说明 \(\sigma(A)\) 有界.

现在设 \(\lambda_{0}\in\rho(A)\).对另一个 \(\lambda\in\mathbb{C}\),有
\[
A-\lambda I=(A-\lambda_{0}I)\left(I-(\lambda-\lambda_{0})(A-\lambda_{0}I)^{-1}\right).
\label{eq:resolvent-local-factorization}
\]
若
\[
|\lambda-\lambda_{0}|<\frac{1}{\|(A-\lambda_{0}I)^{-1}\|},
\]
则式~\eqref{eq:resolvent-local-factorization} 右端的两个因子都可逆,所以 \(A-\lambda I\) 也可逆.这说明预解集是开集,谱是闭集.

最后证明谱非空.由式~\eqref{eq:resolvent-local-factorization} 可得,在 \(\lambda_{0}\) 附近
\[
\begin{aligned}
(A-\lambda I)^{-1}
&=\left(I-(\lambda-\lambda_{0})(A-\lambda_{0}I)^{-1}\right)^{-1}(A-\lambda_{0}I)^{-1}\\
&=\sum_{m=0}^{\infty}(\lambda-\lambda_{0})^{m}
\bigl((A-\lambda_{0}I)^{-1}\bigr)^{m}(A-\lambda_{0}I)^{-1}.
\end{aligned}
\label{eq:resolvent-local-series}
\]
因此,对任意 \(\phi,\psi\in\mathcal{H}\),函数
\[
\lambda\longmapsto\langle\phi,(A-\lambda I)^{-1}\psi\rangle
\label{eq:resolvent-holomorphic}
\]
在预解集上全纯;并且由式~\eqref{eq:resolvent-neumann},当 \(|\lambda|\to\infty\) 时它趋于零.

若 \(\sigma(A)\) 为空,则式~\eqref{eq:resolvent-holomorphic} 在整个 \(\mathbb{C}\) 上全纯且在无穷远处趋于零.根据 Liouville 定理,它恒为零,对所有 \(\phi,\psi\) 都如此,因而 \((A-\lambda I)^{-1}\) 是零算子;但这不可能,因为
\[
(A-\lambda I)(A-\lambda I)^{-1}=I.
\]
故谱非空.
\end{proof}

注意,若 \(A\psi=\lambda\psi\),其中 \(\psi\neq0\),则 \(A-\lambda I\) 有非零核,所以 \(\lambda\in\sigma(A)\).因此,任何特征值都属于谱.不过在无限维情形中,反过来未必成立:谱中的点可能不是特征值.对于有界自伴算子,谱仍可以用一种与有限维情形相当接近的方式刻画.

\begin{proposition}[有界自伴算子的谱刻画]\label{prop:selfadjoint-spectrum-characterization}
若 \(A\in\mathcal{B}(\mathcal{H})\) 自伴,则:
\begin{enumerate}
\item \(\sigma(A)\subset\mathbb{R}\);
\item 对 \(\lambda\in\mathbb{R}\),\(\lambda\in\sigma(A)\) 当且仅当存在一列 \(\mathcal{H}\) 中的非零向量 \(\psi_{n}\),使得
\[
\lim_{n\to\infty}\frac{\|A\psi_{n}-\lambda\psi_{n}\|}{\|\psi_{n}\|}=0.
\label{eq:almost-eigenvalue}
\]
\end{enumerate}
\end{proposition}

命题中的第二点说,\(\lambda\) 属于谱,当且仅当它是一个“几乎特征值”:存在非零 \(\psi\),使得 \(A\psi\) 等于 \(\lambda\psi\) 加上一个相对于 \(\psi\) 的大小可以任意小的误差.

\begin{lemma}[自伴算子的下界估计]\label{lem:selfadjoint-resolvent-estimate}
若 \(A\in\mathcal{B}(\mathcal{H})\) 自伴,且 \(\lambda=a+ib\in\mathbb{C}\),则
\[
\langle(A-\lambda I)\psi,(A-\lambda I)\psi\rangle
\geq b^{2}\langle\psi,\psi\rangle.
\label{eq:selfadjoint-resolvent-estimate}
\]
\end{lemma}

\begin{proof}
直接计算得
\begin{align*}
&\langle(A-(a+ib)I)\psi,(A-(a+ib)I)\psi\rangle\\
&=\langle(A-aI)\psi,(A-aI)\psi\rangle
 +ib\langle\psi,(A-aI)\psi\rangle
 -ib\langle(A-aI)\psi,\psi\rangle
 +b^{2}\langle\psi,\psi\rangle.
\end{align*}
由于 \(A-aI\) 自伴,右端第二项和第三项相消.因此
\[
\langle(A-\lambda I)\psi,(A-\lambda I)\psi\rangle
=\langle(A-aI)\psi,(A-aI)\psi\rangle+b^{2}\langle\psi,\psi\rangle,
\]
从而得到所需不等式.
\end{proof}

\begin{proof}[命题~\ref{prop:selfadjoint-spectrum-characterization} 的证明]
先证第一点.设 \(\lambda=a+ib\),且 \(b\neq0\).由引理~\ref{lem:selfadjoint-resolvent-estimate},\(A-\lambda I\) 是单射.由命题~\ref{prop:range-kernel-adjoint},
\[
\operatorname{Range}(A-\lambda I)^{\perp}=\ker(A-\overline{\lambda}I).
\]
因为 \(\overline{\lambda}\) 也有非零虚部,\(A-\overline{\lambda}I\) 是单射,故 \(A-\lambda I\) 的值域在 \(\mathcal{H}\) 中稠密.取任意 \(\phi\in\mathcal{H}\),令 \(\phi_{n}=(A-\lambda I)\psi_{n}\) 是值域中的向量列,且 \(\phi_{n}\to\phi\).对 \(\psi_{n}-\psi_{m}\) 使用引理中的估计,得到 \(\{\psi_{n}\}\) 是 Cauchy 列.因此 \(\psi_{n}\to\psi\);由于 \(A\) 有界,
\[
(A-\lambda I)\psi=\lim_{n\to\infty}(A-\lambda I)\psi_{n}=\lim_{n\to\infty}\phi_{n}=\phi.
\]
故 \(A-\lambda I\) 一一且满射.上述估计也说明其逆有界(也可使用闭图定理),所以 \(\lambda\in\rho(A)\).于是谱包含于实轴.

再证第二点.若存在式~\eqref{eq:almost-eigenvalue} 中的向量列,且 \(A-\lambda I\) 有逆,令 \(\phi_{n}=(A-\lambda I)\psi_{n}\),则 \(\psi_{n}=(A-\lambda I)^{-1}\phi_{n}\),从而
\[
\lim_{n\to\infty}\frac{\|\phi_{n}\|}{\|(A-\lambda I)^{-1}\phi_{n}\|}=0,
\]
这说明 \((A-\lambda I)^{-1}\) 实际上无界,矛盾.因此 \(A-\lambda I\) 没有有界逆.

反过来,若不存在满足式~\eqref{eq:almost-eigenvalue} 的向量列,则存在 \(\varepsilon>0\),使得对所有 \(\psi\in\mathcal{H}\),
\[
\|(A-\lambda I)\psi\|\geq\varepsilon\|\psi\|.
\label{eq:selfadjoint-lower-bound}
\]
于是 \(A-\lambda I\) 是单射;并且它是自伴算子,由命题~\ref{prop:range-kernel-adjoint},其值域稠密.利用式~\eqref{eq:selfadjoint-lower-bound},像前面的论证一样可知其值域也是闭集,因而等于整个 \(\mathcal{H}\).所以 \(A-\lambda I\) 有有界逆.
\end{proof}

\begin{example}[乘法算子的谱]\label{ex:multiplication-spectrum}
令 \(\mathcal{H}=L^{2}([0,1])\),定义
\[
(A\psi)(x)=x\psi(x).
\]
则 \(A\) 有界且自伴,并且
\[
\sigma(A)=[0,1].
\]
如第~\ref{sec:spectral-infinite-difficulties} 节所述,\(A\) 没有真正的特征向量.
\end{example}

\begin{proof}
显然 \(\|A\psi\|\leq\|\psi\|\),并且对所有 \(\phi,\psi\in\mathcal{H}\),有 \(\langle\phi,A\psi\rangle=\langle A\phi,\psi\rangle\),故 \(A\) 有界且自伴.给定 \(\lambda\in(0,1)\),考虑
\[
\psi_{n}:=\mathbf{1}_{[\lambda,\lambda+1/n]}.
\]
当 \(n\) 足够大时,\(\|\psi_{n}\|^{2}=1/n\).另一方面,在区间 \([\lambda,\lambda+1/n]\) 上有 \(|x-\lambda|\leq1/n\),所以
\[
\|(A-\lambda I)\psi_{n}\|^{2}\leq\frac{1}{n^{3}}.
\]
由命题~\ref{prop:selfadjoint-spectrum-characterization},\(\lambda\in\sigma(A)\).这对所有 \(\lambda\in(0,1)\) 都成立,而谱是闭集,故 \(\sigma(A)\supset[0,1]\).反之,若 \(\lambda\notin[0,1]\),则函数 \(1/(x-\lambda)\) 在 \([0,1]\) 上有界,\(A-\lambda I\) 的逆就是乘以该函数的算子.因此 \(\sigma(A)=[0,1]\).
\end{proof}

\section{有界自伴算子的谱定理(一)}
\label{sec:spectral-theorem-first}

\subsection{谱子空间}
\label{subsec:spectral-subspaces}

给定一个有界(暂时如此)自伴算子 \(A\),我们希望把每个 Borel 集 \(E\subset\sigma(A)\) 关联到 \(\mathcal{H}\) 的一个闭子空间 \(V_{E}\).直观上,\(V_{E}\) 应当是特征值属于 \(E\) 的广义特征向量所张成的闭子空间.我们希望这些子空间满足:
\begin{enumerate}
\item \(V_{\sigma(A)}=\mathcal{H}\),且 \(V_{\varnothing}=\{0\}\);
\item 若 \(E,F\) 不相交,则 \(V_{E}\perp V_{F}\);
\item 对任意 \(E,F\),有 \(V_{E\cap F}=V_{E}\cap V_{F}\);
\item 若 \(E_{1},E_{2},\ldots\) 两两不相交,且 \(E=\bigcup_{j}E_{j}\),则
\[
V_{E}=\bigoplus_{j}V_{E_{j}};
\]
\item 对任意 \(E\),\(V_{E}\) 在 \(A\) 下不变;
\item 若 \(E\subset[\lambda_{0}-\varepsilon,\lambda_{0}+\varepsilon]\),且 \(\psi\in V_{E}\),则
\[
\|(A-\lambda_{0}I)\psi\|\leq\varepsilon\|\psi\|.
\]
\end{enumerate}
条件 \(V_{\sigma(A)}=\mathcal{H}\) 表达了广义特征向量应张成整个 \(\mathcal{H}\);条件 2 表达了不同特征值对应的广义特征向量仍应具有某种正交性,即使它们并不真正属于 Hilbert 空间.条件 4 中如果出现无穷多个集合,直和取 Hilbert 空间意义下的直和(定义~\ref{app:def:hilbert-direct-sum}).条件 5 和条件 6 则表达了 \(V_{E}\) 由特征值位于 \(E\) 中的广义特征向量构成.

\subsection{投影值测度}
\label{subsec:projection-valued-measures}

用相应的正交投影算子描述 Hilbert 空间的闭子空间很方便.回忆附录~\ref{app:prop:orthogonal-projection}:给定 \(\mathcal{H}\) 的闭子空间 \(V\),存在唯一有界算子 \(P\),它在 \(V\) 上等于恒等算子,在正交补 \(V^{\perp}\) 上等于零.该算子称为投影到 \(V\) 上的正交投影,并满足 \(P^{2}=P\) 及 \(P^{*}=P\).下面的定义把谱子空间的前四个性质(即不涉及算子 \(A\) 的性质)表述为相应正交投影的性质.由于这些性质与测度的性质相似,我们称其为投影值测度.

\begin{definition}[投影值测度]\label{def:projection-valued-measure}
设 \(X\) 是集合,\(\Omega\) 是 \(X\) 上的 \(\sigma\)-代数.若映射 \(\mu:\Omega\to\mathcal{B}(\mathcal{H})\) 满足下列性质,则称其为投影值测度:
\begin{enumerate}
\item 对每个 \(E\in\Omega\),\(\mu(E)\) 都是正交投影;
\item \(\mu(\varnothing)=0\),且 \(\mu(X)=I\);
\item 若 \(E_{1},E_{2},\ldots\in\Omega\) 两两不相交,则对所有 \(v\in\mathcal{H}\),有
\[
\mu\left(\bigcup_{j=1}^{\infty}E_{j}\right)v
=\sum_{j=1}^{\infty}\mu(E_{j})v,
\]
其中级数在 \(\mathcal{H}\) 的范数拓扑中收敛;
\item 对所有 \(E_{1},E_{2}\in\Omega\),有
\[
\mu(E_{1}\cap E_{2})=\mu(E_{1})\mu(E_{2}).
\]
\end{enumerate}
\end{definition}

若 \(E_{1},E_{2}\) 不相交,则定义中的性质 2 和性质 4 告诉我们 \(\mu(E_{1})\mu(E_{2})=0\),由此可知 \(\mu(E_{1})\) 与 \(\mu(E_{2})\) 的值域相互正交(见练习~\ref{ex:pvm-orthogonality}).不难验证,\(\mu(E_{1})\mu(E_{2})\) 是投影到两个值域交集上的投影(见练习~\ref{ex:pvm-intersection}).因此,若对每个 \(E\in\Omega\) 定义闭子空间 \(V_{E}:=\operatorname{Range}(\mu(E))\),则这些 \(V_{E}\) 满足我们对谱子空间预期的前四个性质.

下一小节中,我们将给每个有界自伴算子 \(A\) 关联一个投影值测度 \(\mu_{A}\).此时,投影 \(\mu_{A}(E)\) 可以理解为对应于 \(E\) 的谱子空间上的投影.我们还将引入关于投影值测度的算子值积分;对于与 \(A\) 关联的投影值测度,这个算子值积分正是 \(A\) 的函数演算.

注意,对任意投影值测度 \(\mu\) 和 \(\psi\in\mathcal{H}\),可定义普通的正实值测度 \(\mu_{\psi}\):
\[
\mu_{\psi}(E)=\langle\psi,\mu(E)\psi\rangle.
\label{eq:pvm-scalar-measure}
\]
这就建立了关于投影值测度积分与关于普通测度积分之间的联系.

\begin{proposition}[算子值积分]\label{prop:operator-valued-integration}
设 \(\Omega\) 是集合 \(X\) 上的 \(\sigma\)-代数,\(\mu:\Omega\to\mathcal{B}(\mathcal{H})\) 是投影值测度.则存在唯一的线性映射
\[
f\longmapsto\int_{X}f\,\mathrm{d}\mu
\]
从有界可测复值函数空间映到 \(\mathcal{B}(\mathcal{H})\),使得对所有这样的 \(f\) 以及所有 \(\psi\in\mathcal{H}\),
\[
\left\langle\psi,\left(\int_{X}f\,\mathrm{d}\mu\right)\psi\right\rangle
=\int_{X}f\,\mathrm{d}\mu_{\psi}.
\label{eq:operator-valued-integral-definition}
\]
此外:
\begin{enumerate}
\item 对每个 \(E\in\Omega\),有
\[
\int_{X}\mathbf{1}_{E}\,\mathrm{d}\mu=\mu(E),
\]
特别地,常值函数 1 的积分是 \(I\);
\item 对所有 \(f\),有
\[
\left\|\int_{X}f\,\mathrm{d}\mu\right\|
\leq\sup_{\lambda\in X}|f(\lambda)|;
\label{eq:operator-valued-integral-bound}
\]
\item 积分具有乘法性:
\[
\int_{X}fg\,\mathrm{d}\mu
=\left(\int_{X}f\,\mathrm{d}\mu\right)
 \left(\int_{X}g\,\mathrm{d}\mu\right);
\label{eq:operator-valued-integral-multiplicative}
\]
\item 对所有 \(f\),有
\[
\int_{X}\overline{f}\,\mathrm{d}\mu
=\left(\int_{X}f\,\mathrm{d}\mu\right)^{*}.
\]
特别地,若 \(f\) 为实值函数,则 \(\int_{X}f\,\mathrm{d}\mu\) 自伴.
\end{enumerate}
\end{proposition}

由性质 1 和线性性,投影值测度积分对简单函数具有预期行为.再由性质 2,任意有界可测函数 \(f\) 的积分都可这样计算:取一致收敛于 \(f\) 的简单函数列 \(s_{n}\),则 \(f\) 的积分就是 \(s_{n}\) 的积分在算子范数拓扑中的极限.

积分的乘法性初看起来也许令人惊讶.不过,对任意 \(E_{1},E_{2}\in\Omega\),定义中的性质 3 给出
\[
\left(\int_{X}\mathbf{1}_{E_{1}}\,\mathrm{d}\mu\right)
\left(\int_{X}\mathbf{1}_{E_{2}}\,\mathrm{d}\mu\right)
=\mu(E_{1})\mu(E_{2})=\mu(E_{1}\cap E_{2})
=\int_{X}\mathbf{1}_{E_{1}}\mathbf{1}_{E_{2}}\,\mathrm{d}\mu.
\]
所以,示性函数层面的乘法性已经包含在投影值测度的定义中.

如果试图构造一种实值测度,使相应积分具有乘法性,那么由于 \(\mathbf{1}_{E}^{2}=\mathbf{1}_{E}\),\(\mathbf{1}_{E}\) 的积分,即 \(\mu(E)\),必须满足 \(\mu(E)^{2}=\mu(E)\).这意味着对每个 \(E\),\(\mu(E)\) 只能是 0 或 1.这样的测度确实具有积分的乘法性,但并不有趣.对于算子值测度,\(\mathcal{B}(\mathcal{H})\) 中幂等元(满足 \(A^{2}=A\) 的元素)远多于 \(\mathbb{R}\) 中的幂等元,因此可以得到有趣的乘法性例子.

\begin{proof}[命题~\ref{prop:operator-valued-integration} 的证明]
给定投影值测度 \(\mu\) 和有界可测函数 \(f\),定义 \(Q_{f}:\mathcal{H}\to\mathbb{C}\):
\[
Q_{f}(\psi)=\int_{X}f\,\mathrm{d}\mu_{\psi},
\]
其中 \(\mu_{\psi}\) 由式~\eqref{eq:pvm-scalar-measure} 给出.若 \(f\) 是示性函数,则 \(Q_{f}(\psi)=\langle\psi,\mu(E)\psi\rangle\) 是有界二次型(见定义~\ref{app:def:quadratic-form}).从示性函数到简单函数,再到一般有界可测函数的推广是直接的:对任何有界可测 \(f\),\(Q_{f}\) 都是有界二次型,并满足
\[
|Q_{f}(\psi)|\leq\sup_{\lambda\in X}|f(\lambda)|\,\|\psi\|^{2}.
\label{eq:operator-valued-quadratic-bound}
\]
由命题~\ref{app:prop:quadratic-operator},存在唯一有界算子 \(A_{f}\),使得
\[
Q_{f}(\psi)=\langle\psi,A_{f}\psi\rangle.
\]
定义 \(\int_{X}f\,\mathrm{d}\mu:=A_{f}\),便得到式~\eqref{eq:operator-valued-integral-definition}.线性映射的唯一性来自同一命题中的唯一性.

若 \(f=\mathbf{1}_{E}\),则 \(Q_{f}(\psi)=\mu_{\psi}(E)=\langle\psi,\mu(E)\psi\rangle\),所以相应算子 \(A_{f}=\mu(E)\),这证明了性质 1.性质 2 由式~\eqref{eq:operator-valued-quadratic-bound} 得到.

对于性质 3,我们已经观察到,示性函数层面的乘法性由投影值测度的定义保证.由于式~\eqref{eq:operator-valued-integral-multiplicative} 两边关于 \((\phi,\psi)\) 都是双线性的,它对简单函数成立;再使用性质 2 取极限,即可推广到所有有界可测函数.

最后,若 \(f\) 为实值函数,则 \(Q_{f}(\psi)\) 对所有 \(\psi\) 都是实数.由命题~\ref{app:prop:quadratic-operator},相应算子 \(A_{f}\) 自伴,性质 4 随即由线性性得到.
\end{proof}

\subsection{谱定理}
\label{subsec:spectral-theorem-first-form}

现在可以陈述有界自伴算子谱定理的一种形式.

\begin{theorem}[谱定理:第一种形式]\label{thm:spectral-theorem-first}
若 \(A\in\mathcal{B}(\mathcal{H})\) 自伴,则在 \(\sigma(A)\) 上的 Borel \(\sigma\)-代数存在唯一的投影值测度 \(\mu_{A}\),其值为 \(\mathcal{H}\) 上的投影,并满足
\[
\int_{\sigma(A)}\lambda\,\mathrm{d}\mu_{A}(\lambda)=A.
\label{eq:spectral-theorem-first}
\]
\end{theorem}

由于 \(A\) 的谱有界,函数 \(f(\lambda)=\lambda\) 在 \(\sigma(A)\) 上有界.本定理的证明将在第~8 章给出.

\begin{definition}[函数演算]\label{def:functional-calculus}
设 \(A\in\mathcal{B}(\mathcal{H})\) 自伴,\(f:\sigma(A)\to\mathbb{C}\) 是有界可测函数.令
\[
f(A):=\int_{\sigma(A)}f(\lambda)\,\mathrm{d}\mu_{A}(\lambda),
\label{eq:functional-calculus}
\]
其中 \(\mu_{A}\) 是定理~\ref{thm:spectral-theorem-first} 中的投影值测度,则称 \(f(A)\) 为 \(f\) 在 \(A\) 处的函数演算.
\end{definition}

把 \(\mu_{A}\) 从 \(\sigma(A)\) 延拓到整个 \(\mathbb{R}\),规定 \(\mathbb{R}\setminus\sigma(A)\) 的测度为零.粗略地说,\(f(A)\) 在投影算子 \(\mu_{A}([\lambda,\lambda+\mathrm{d}\lambda))\) 的值域上等于 \(f(\lambda)I\).

由于关于 \(\mu_{A}\) 的积分具有乘法性,由式~\eqref{eq:spectral-theorem-first} 可知,如果 \(f(\lambda)=\lambda^{m}\),其中 \(m\) 是正整数,则 \(f(A)=A^{m}\).进一步地,级数 \(e^{a\lambda}=\sum_{m=0}^{\infty}(a\lambda)^{m}/m!\) 在紧集 \(\sigma(A)\) 上一致收敛,所以由函数 \(f(\lambda)=e^{a\lambda}\) 的函数演算得到的 \(e^{aA}\) 可以用幂级数计算.

\begin{definition}[谱子空间]\label{def:spectral-subspace}
设 \(A\in\mathcal{B}(\mathcal{H})\) 自伴,令 \(\mu_{A}\) 为相应的投影值测度,并将它延拓到 \(\mathbb{R}\),令 \(\mathbb{R}\setminus\sigma(A)\) 的测度为零.对每个 Borel 集 \(E\subset\mathbb{R}\),定义 \(\mathcal{H}\) 的谱子空间为
\[
V_{E}:=\operatorname{Range}(\mu_{A}(E)).
\label{eq:spectral-subspace-definition}
\]
\end{definition}

投影值测度的定义保证这些谱子空间满足第~\ref{subsec:spectral-subspaces} 节列出的前四个性质.下面说明式~\eqref{eq:spectral-theorem-first} 如何给出剩下的两个性质.

\begin{proposition}[谱子空间的性质]\label{prop:spectral-subspace-properties}
若 \(A\in\mathcal{B}(\mathcal{H})\) 自伴,则由 \(A\) 关联的谱子空间满足:
\begin{enumerate}
\item 每个谱子空间 \(V_{E}\) 在 \(A\) 下不变;
\item 若 \(E\subset[\lambda_{0}-\varepsilon,\lambda_{0}+\varepsilon]\),则对所有 \(\psi\in V_{E}\),有
\[
\|(A-\lambda_{0}I)\psi\|\leq\varepsilon\|\psi\|;
\]
\item \(A|_{V_{E}}\) 的谱包含于 \(E\) 的闭包;
\item 若 \(\lambda_{0}\in\sigma(A)\),则对 \(\lambda_{0}\) 的每个邻域 \(U\),都有 \(V_{U}\neq\{0\}\),等价地,\(\mu_{A}(U)\neq0\).
\end{enumerate}
\end{proposition}

\begin{proof}
先证第一点.对 \(\sigma(A)\) 上任意两个有界可测函数 \(f,g\),算子 \(f(A)\) 和 \(g(A)\) 可交换,因为它们按任意次序相乘都等于函数 \(fg=gf\) 的积分.特别地,\(A\)(即函数 \(f(\lambda)=\lambda\) 的积分)与 \(\mu_{A}(E)\)(即示性函数 \(\mathbf{1}_{E}\) 的积分)可交换.因此,对值域中的向量 \(\mu_{A}(E)\phi\),有
\[
A\mu_{A}(E)\phi=\mu_{A}(E)A\phi,
\]
右端仍在 \(\mu_{A}(E)\) 的值域中,故谱子空间不变.

对于第二点,设 \(\psi\in V_{E}\),且 \(E\subset[\lambda_{0}-\varepsilon,\lambda_{0}+\varepsilon]\).则 \(\psi\) 在 \(\mu_{A}(E)\) 的值域中,于是
\[
(A-\lambda_{0}I)\psi=(A-\lambda_{0}I)\mu_{A}(E)\psi.
\]
令 \(f(\lambda)=\lambda-\lambda_{0}\),则 \(\mu_{A}(E)=\mathbf{1}_{E}(A)\),且 \(A-\lambda_{0}I=f(A)\).由积分的乘法性,
\[
(A-\lambda_{0}I)\psi=(f\mathbf{1}_{E})(A)\psi.
\]
因为 \(|f(\lambda)\mathbf{1}_{E}(\lambda)|\leq\varepsilon\),由式~\eqref{eq:operator-valued-integral-bound},算子 \((f\mathbf{1}_{E})(A)\) 的范数至多为 \(\varepsilon\).

对于第三点,若 \(\lambda_{0}\notin\overline{E}\),则函数
\[
g(\lambda)=\frac{\mathbf{1}_{E}(\lambda)}{\lambda-\lambda_{0}}
\]
有界.于是 \(g(A)\) 有界,并且
\[
g(A)(A-\lambda_{0}I)=(A-\lambda_{0}I)g(A)=\mathbf{1}_{E}(A).
\]
这说明 \(g(A)\) 在 \(V_{E}\) 上的限制是 \(A\) 在 \(V_{E}\) 上的限制的逆.因此 \(\lambda_{0}\notin\sigma(A|_{V_{E}})\).

对于第四点,固定 \(\lambda_{0}\in\sigma(A)\).假设对某个 \(\varepsilon>0\),有 \(\mu_{A}((\lambda_{0}-\varepsilon,\lambda_{0}+\varepsilon))=0\).考虑有界函数
\[
f(\lambda)=
\begin{cases}
\dfrac{1}{\lambda-\lambda_{0}},&|\lambda-\lambda_{0}|\geq\varepsilon,\\
0,&|\lambda-\lambda_{0}|<\varepsilon.
\end{cases}
\]
函数 \(f(\lambda)(\lambda-\lambda_{0})\) 除上述区间外恒等于 1;由于该区间的测度为零,它与常值函数 1 在 \(\mu_{A}\)-几乎处处意义下相同.因此由积分的乘法性,
\[
f(A)(A-\lambda_{0}I)=(A-\lambda_{0}I)f(A)=I,
\]
说明 \(f(A)\) 是 \(A-\lambda_{0}I\) 的有界逆,矛盾于 \(\lambda_{0}\in\sigma(A)\).
\end{proof}

\begin{proposition}[与算子相交换的算子]\label{prop:commuting-operator-functional-calculus}
若 \(A\in\mathcal{B}(\mathcal{H})\) 自伴,且 \(B\in\mathcal{B}(\mathcal{H})\) 与 \(A\) 相交换,则:
\begin{enumerate}
\item 对 \(\sigma(A)\) 上任意有界可测函数 \(f\),\(f(A)\) 与 \(B\) 相交换;
\item \(A\) 的每个谱子空间在 \(B\) 下不变.
\end{enumerate}
\end{proposition}

本命题的证明留到第~8 章.到这里,我们至少对有界自伴算子实现了谱定理的一个目标:给出描述自伴算子 \(A\) 在状态 \(\psi\) 中测量概率的概率测度.

\begin{proposition}[测量概率的谱测度]\label{prop:spectral-measurement-probability}
设 \(A\in\mathcal{B}(\mathcal{H})\) 自伴,\(\psi\in\mathcal{H}\) 是单位向量.则存在唯一的 \(\mathbb{R}\) 上概率测度 \(\mu_{\psi}^{A}\),使得对所有非负整数 \(m\),
\[
\int_{\mathbb{R}}\lambda^{m}\,\mathrm{d}\mu_{\psi}^{A}(\lambda)
=\langle\psi,A^{m}\psi\rangle.
\label{eq:moment-spectral-measure}
\]
\end{proposition}

第~9 章将证明无界自伴算子版本的本命题.不过在无界情形,即使 \(\psi\) 属于所有 \(A^{m}\) 的定义域,我们也无法得到概率测度的唯一性;谱定理仍提供一个规范的选择.

\begin{proof}
按第~\ref{subsec:projection-valued-measures} 节的方式,在 \(\sigma(A)\) 上定义
\[
\mu_{\psi}^{A}(E)=\langle\psi,\mu_{A}(E)\psi\rangle.
\]
投影值测度积分的性质给出
\[
\langle\psi,A^{m}\psi\rangle
=\left\langle\psi,\left(\int_{\sigma(A)}\lambda^{m}\,\mathrm{d}\mu_{A}(\lambda)\right)\psi\right\rangle
=\int_{\sigma(A)}\lambda^{m}\,\mathrm{d}\mu_{\psi}^{A}(\lambda).
\]
把 \(\mu_{\psi}^{A}\) 延拓到 \(\mathbb{R}\),在 \(\mathbb{R}\setminus\sigma(A)\) 上令其为零,即得到所需概率测度.又因为
\[
|\langle\psi,A^{m}\psi\rangle|\leq\|\psi\|^{2}\|A^{m}\|\leq\|\psi\|^{2}\|A\|^{m},
\]
其矩至多指数增长,关于矩问题的标准唯一性定理(例如文献~[18] 第 4 章定理 8.1)给出 \(\mu_{\psi}^{A}\) 的唯一性.
\end{proof}

\ChapterSevenLater
